\begin{figure}[H]
  \centering
  \begin{tikzpicture}
    \begin{axis}[paperaxis,
      xlabel={Année}, ylabel={Écart France moins comparables (part de l'actif)},
      xtick={2010,2011,2012,2013,2014,2015,2016,2017,2018,2019,2020,2021,2022},
      xticklabel style={/pgf/number format/1000 sep=, rotate=45, anchor=east},
      legend pos=north west, legend cell align=left,
    ]
      \draw[zeroline, line width=0.6pt] (axis cs:2010,0) -- (axis cs:2022,0);
      % Deux marges du capital retenu de la France, toutes deux en bleu (France) ; le style de trait distingue.
      \addplot[franceblue, line width=1.3pt, mark=*, mark options={scale=0.7}]
        table[col sep=comma, x=year, y=retention] {figures/capital_retenu.csv};
      \addlegendentry{Réserves sur actif}
      \addplot[franceblue, line width=1.3pt, densely dashed, mark=square*, mark options={scale=0.6}]
        table[col sep=comma, x=year, y=equity] {figures/capital_retenu.csv};
      \addlegendentry{Fonds propres sur actif}
      % Les deux chocs, chacun dans sa couleur : orange la hausse de 2013, vert la baisse de 2018.
      \draw[hausse2013, dashed, line width=1pt]
        ({axis cs:2012.5,0}|-{rel axis cs:0,0}) -- ({axis cs:2012.5,0}|-{rel axis cs:0,1});
      \draw[baisse2018, dashed, line width=1pt]
        ({axis cs:2017.5,0}|-{rel axis cs:0,0}) -- ({axis cs:2017.5,0}|-{rel axis cs:0,1});
    \end{axis}
  \end{tikzpicture}
  \caption{Niveau de capital retenu de la France, référence 2012, autour des deux chocs}
  \label{fig:capital-retenu}
  \fignote{Coefficients $\tau_j$ de l'équation~\eqref{eq:event} sur le niveau de capital retenu,
  réserves sur actif et fonds propres sur actif, écart de la France aux cinq pays comparables l'année $j$,
  année de référence 2012 omise, panel étendu à 2010. Les deux séries sont en bleu, la couleur de la
  France dans tout le papier ; le trait plein est la rétention, le trait tireté les fonds propres.
  La ligne orange verticale marque la hausse d'impôt de 2013, la ligne verte la baisse de 2018. Un
  profil plat avant 2013, un saut en 2013, puis un plateau sans retour en 2018 est la forme qui
  soutient le cliquet.}
\end{figure}
